\chapter{Optimization in Practice}

\section{Stochastic Gradient Descent}

In Chapter 4, we introduced gradient-based optimization as the primary method for training machine learning models. In practice, however, computing full gradients over large datasets is often computationally infeasible.

\medskip

\noindent
\textbf{Stochastic Gradient Descent (SGD).}  
Instead of computing the full gradient,
\[
\nabla L(\theta) = \frac{1}{n} \sum_{i=1}^n \nabla \ell(f_\theta(x_i), y_i),
\]
we approximate it using a single randomly selected data point:
\[
\theta_{t+1} = \theta_t - \eta \nabla \ell(f_\theta(x_{i_t}), y_{i_t}).
\]

\medskip

\noindent
\textbf{Interpretation.}
\begin{itemize}
    \item Each update is a noisy estimate of the true gradient
    \item The noise arises from sampling
\end{itemize}

\medskip

\noindent
\textbf{Key properties.}
\begin{itemize}
    \item Low computational cost per iteration
    \item Enables training on large datasets
    \item Introduces stochasticity into optimization
\end{itemize}

\medskip

\noindent
\textbf{Convergence behavior (informal).}
\begin{itemize}
    \item SGD does not converge exactly to a minimum
    \item Instead, it fluctuates around regions of low loss
\end{itemize}

\medskip

\noindent
\textbf{Insight.}  
The stochasticity of SGD can help escape sharp minima and explore the loss landscape more effectively.

\section{Mini-Batching and Efficiency}

In practice, SGD is rarely implemented using single data points. Instead, we use mini-batches.

\medskip

\noindent
\textbf{Mini-batch SGD.}
\[
\theta_{t+1} = \theta_t - \eta \frac{1}{|B_t|} \sum_{i \in B_t} \nabla \ell(f_\theta(x_i), y_i),
\]
where $B_t$ is a subset (batch) of data points.

\medskip

\noindent
\textbf{Advantages.}
\begin{itemize}
    \item Reduces variance of gradient estimates
    \item Exploits parallel computation (e.g., GPUs)
    \item More stable updates than pure SGD
\end{itemize}

\medskip

\noindent
\textbf{Batch size tradeoff.}
\begin{itemize}
    \item Small batch:
    \begin{itemize}
        \item Noisy updates
        \item Better generalization in some cases
    \end{itemize}
    \item Large batch:
    \begin{itemize}
        \item More accurate gradients
        \item Higher computational cost per step
    \end{itemize}
\end{itemize}

\medskip

\noindent
\textbf{Efficiency considerations.}
\begin{itemize}
    \item Memory constraints limit batch size
    \item Larger batches benefit from hardware acceleration
\end{itemize}

\medskip

\noindent
\textbf{Epochs and iterations.}
\begin{itemize}
    \item One epoch: one pass through the dataset
    \item Iteration: one parameter update
\end{itemize}

\section{Momentum and Adaptive Methods}

Basic SGD can be slow or unstable, especially in high-dimensional problems. Several enhancements improve its performance.

\subsection{Momentum}

\medskip

\noindent
\textbf{Update rule.}
\[
v_{t+1} = \beta v_t + \nabla L(\theta_t), \quad
\theta_{t+1} = \theta_t - \eta v_{t+1}.
\]

\medskip

\noindent
\textbf{Interpretation.}
\begin{itemize}
    \item Accumulates past gradients
    \item Smooths updates and accelerates convergence
\end{itemize}

\medskip

\noindent
\textbf{Geometric intuition.}  
Momentum helps traverse flat regions quickly and reduces oscillations in directions of high curvature.

\subsection{Adaptive Methods}

Adaptive methods adjust learning rates for individual parameters.

\medskip

\noindent
\textbf{Examples.}
\begin{itemize}
    \item AdaGrad
    \item RMSProp
    \item Adam
\end{itemize}

\medskip

\noindent
\textbf{Key idea.}
\begin{itemize}
    \item Scale updates based on past gradient magnitudes
    \item Use larger steps for infrequent features
    \item Use smaller steps for frequently updated parameters
\end{itemize}

\medskip

\noindent
\textbf{Adam (conceptual form).}
\begin{itemize}
    \item Maintains moving averages of gradients and squared gradients
    \item Combines momentum and adaptive scaling
\end{itemize}

\medskip

\noindent
\textbf{Practical observation.}
\begin{itemize}
    \item Adaptive methods often converge faster
    \item SGD with momentum may generalize better in some settings
\end{itemize}

\section{Practical Convergence Behavior}

In real-world settings, optimization rarely follows idealized theoretical behavior.

\medskip

\noindent
\textbf{Learning rate scheduling.}
\begin{itemize}
    \item Decrease learning rate over time
    \item Common strategies:
    \begin{itemize}
        \item step decay,
        \item exponential decay,
        \item cosine annealing
    \end{itemize}
\end{itemize}

\medskip

\noindent
\textbf{Warm-up.}
\begin{itemize}
    \item Start with small learning rates
    \item Gradually increase to stabilize early training
\end{itemize}

\medskip

\noindent
\textbf{Plateaus and sharp drops.}
\begin{itemize}
    \item Training loss may decrease slowly, then suddenly improve
    \item Reflects transitions in the loss landscape
\end{itemize}

\medskip

\noindent
\textbf{Noise and generalization.}
\begin{itemize}
    \item Noise from SGD can act as implicit regularization
    \item Helps avoid sharp minima
\end{itemize}

\medskip

\noindent
\textbf{Stopping criteria.}
\begin{itemize}
    \item Fixed number of epochs
    \item Convergence of training loss
    \item Validation-based early stopping
\end{itemize}

\medskip

\noindent
\textbf{Practical challenges.}
\begin{itemize}
    \item Choosing learning rates
    \item Tuning hyperparameters
    \item Balancing speed and stability
\end{itemize}

\medskip

\noindent
\textbf{Insight.}  
Optimization in practice is as much an art as it is a science, guided by both theoretical principles and empirical experience.

\section{A Unifying Perspective}

Optimization in modern machine learning involves multiple interacting components:

\begin{itemize}
    \item \textbf{Algorithm:} SGD and its variants
    \item \textbf{Hyperparameters:} learning rate, batch size
    \item \textbf{Dynamics:} noise, curvature, landscape structure
\end{itemize}

\medskip

\noindent
\textbf{Core tradeoffs.}
\begin{itemize}
    \item Speed vs stability
    \item Accuracy vs computational cost
    \item Optimization vs generalization
\end{itemize}

\medskip

\noindent
\textbf{Key insight.}  
The behavior of optimization algorithms shapes not only how fast models learn, but also what solutions they find.

\section{Outlook}

This chapter examined optimization from a practical perspective:
\begin{itemize}
    \item stochastic gradient methods,
    \item mini-batching and efficiency,
    \item momentum and adaptive algorithms,
    \item real-world convergence behavior.
\end{itemize}

\medskip

In the next chapter, we study numerical and statistical stability, focusing on how data preprocessing and conditioning affect learning.

\medskip

\noindent
\textbf{Guiding principle.}  
Effective optimization requires understanding both the algorithm and the structure of the problem it is applied to.